\documentclass[11pt,a4paper]{article}
\usepackage[utf8]{inputenc}
\usepackage[T1]{fontenc}
\usepackage[french]{babel}
\usepackage{lmodern,geometry,microtype}
\usepackage{amsmath}
\usepackage{siunitx}
\usepackage{xcolor,listings,hyperref,booktabs,array}
\geometry{margin=2.4cm}
\sisetup{locale=FR,per-mode=symbol}
\lstset{language=Python,basicstyle=\ttfamily\small,backgroundcolor=\color{black!4},frame=single,breaklines=true,showstringspaces=false}
\hypersetup{colorlinks=true,linkcolor=blue!45!black,urlcolor=blue!55!black}

\title{Documentation de \texttt{generer\_table\_emissivite\_hitran.py}}
\author{Projet climat 2026}
\date{\today}

\begin{document}
\maketitle
\tableofcontents

\section{Rôle du programme}

Ce script est le pilote d'export du modèle. Il calcule l'émissivité spectrale
de chaque couche atmosphérique sur une grille d'altitudes et de longueurs
d'onde, puis écrit le fichier CSV final.

\section{Installation et dépendances}

\begin{lstlisting}[language=bash]
python3.12 -m venv .venv
source .venv/bin/activate
python -m pip install -r requirements.txt
\end{lstlisting}

Le script dépend de toute la chaîne suivante :
\begin{equation*}
\texttt{generer\_table}
\longrightarrow \texttt{code\_final\_emissivite}
\longrightarrow \texttt{2\_sections\_efficaces}
\longrightarrow \text{HAPI et données HITRAN}.
\end{equation*}
Le profil externe
\path{../../temperatures_et_concentrations/code_H20_CH4_CO2_O3.py} est aussi
obligatoire. Une connexion est requise si \path{donnees_hitran/} est incomplet.

\section{Imports du fichier}

\begin{lstlisting}
import csv
from pathlib import Path

import numpy as np

from code_final_emissivite_avec_sections_hitran import (
    calculer_emissivite_une_couche,
    get_gases_ppm,
)
\end{lstlisting}

\texttt{csv} et \texttt{pathlib} font partie de Python. \texttt{numpy} est le
seul paquet externe importé directement. Les deux fonctions locales donnent
respectivement le profil atmosphérique et le calcul radiatif d'une couche.

\section{Colonnes du CSV}

La constante \texttt{COLONNES} fixe 19 champs, dans cet ordre :

\begin{enumerate}
  \item \texttt{altitude\_km}, \texttt{longueur\_onde\_um},
        \texttt{nombre\_onde\_cm\_1} ;
  \item \texttt{pression\_Pa}, \texttt{epaisseur\_couche\_m},
        \texttt{temperature\_K} ;
  \item les concentrations \texttt{CO2/H2O/CH4/O3\_ppm} ;
  \item les quatre sections \texttt{sigma\_*\_m2\_molecule} ;
  \item les quatre contributions \texttt{tau\_*} ;
  \item \texttt{emissivite}.
\end{enumerate}

L'épaisseur optique totale n'est pas une colonne distincte. Elle se retrouve
par la somme des quatre $\tau_g$.

\section{Fonction de génération}

\texttt{generer\_table\_emissivite\_hitran} reçoit :

\begin{center}
\small
\begin{tabular}{lll}
\toprule
Paramètre & Défaut & Signification \\
\midrule
\texttt{fichier\_sortie} & obligatoire & chemin du CSV \\
\texttt{longueur\_onde\_min\_um} & 4,0 & borne basse \\
\texttt{longueur\_onde\_max\_um} & 100,0 & borne haute \\
\texttt{pas\_longueur\_onde\_um} & 0,1 & pas spectral \\
\texttt{altitude\_min\_km} & 0 & altitude minimale \\
\texttt{altitude\_max\_km} & 100 & altitude maximale \\
\texttt{pas\_altitude\_km} & 1 & pas vertical \\
\texttt{epaisseur\_couche\_m} & 1000 & trajet optique vertical \\
\bottomrule
\end{tabular}
\end{center}

La fonction construit les grilles avec \texttt{numpy.arange}, convertit les
longueurs d'onde en mètres et en nombres d'onde, puis ouvre le CSV en UTF-8.
Pour chaque altitude :

\begin{enumerate}
  \item elle obtient pression, température et concentrations ;
  \item elle convertit les degrés Celsius en kelvins ;
  \item elle appelle une seule fois le calcul vectoriel sur toute la grille ;
  \item elle écrit une ligne par longueur d'onde ;
  \item elle vide le tampon sur disque et affiche la progression.
\end{enumerate}

Le retour est le chemin de sortie sous forme d'objet \texttt{Path}.

\section{Dimensions de la sortie par défaut}

La grille contient 101 altitudes et 961 longueurs d'onde, donc
\begin{equation}
101\times961=97061
\end{equation}
lignes de données. Avec l'en-tête, le fichier comporte 97062 lignes. Une
interruption laisse un fichier lisible mais incomplet ; il faut donc contrôler
le nombre de lignes et la dernière altitude.

\section{Exécution}

\begin{lstlisting}[language=bash]
python generer_table_emissivite_hitran.py
\end{lstlisting}

Le fichier produit par défaut est
\path{table_emissivite_hitran_4_100um_0_100km.csv}, dans le dossier du script.
Le calcul est coûteux car quatre spectres Voigt sont recalculés pour chaque
couple de pression et température correspondant à une altitude.

\section{Contrôles recommandés}

Vérifier que $\sigma_g\geq0$, $\tau_g\geq0$ et
$0\leq\varepsilon\leq1$, ainsi que la correspondance
\SI{4}{\micro\metre}$\leftrightarrow$\SI{2500}{\per\centi\metre}. Contrôler
également les 19 noms de colonnes avant d'utiliser le fichier dans un autre
programme.

\end{document}
